% year: תשפ"ב
% type: מבחן
% semester: א
% moed: א
% number: 2ב
% points: 9
% categories: func-analysis
% source: exams/מבחנים/תשפב/מועד א תשפב.pdf

\begin{question}
מצאו תחומי קמירות ונקודות פיתול של הפונקציה
\[ f(x)=x^2\cdot e^{-x} \]
\end{question}

\begin{hint}
חשבו את $f''$ ובדקו את סימנה; שימו לב ש-$e^{-x}>0$ תמיד.
\end{hint}

\begin{solution}
$f$ מוגדרת וגזירה פעמיים בכל $\R$.
\[ f'(x)=2xe^{-x}-x^2e^{-x}=(2x-x^2)e^{-x}, \]
\[ f''(x)=(2-2x)e^{-x}-(2x-x^2)e^{-x}=(x^2-4x+2)e^{-x}. \]
כיוון ש-$e^{-x}>0$, סימן $f''$ הוא סימן $x^2-4x+2$, שאפסיו $x=2\pm\sqrt2$. זהו טרינום עם מקדם מוביל חיובי, ולכן:
\begin{itemize}
  \item $f''>0$ ב-$(-\infty,\,2-\sqrt2)$ וב-$(2+\sqrt2,\,\infty)$: שם $f$ \textbf{קמורה} (קעורה כלפי מעלה).
  \item $f''<0$ ב-$(2-\sqrt2,\,2+\sqrt2)$: שם $f$ \textbf{קעורה} (קעורה כלפי מטה).
\end{itemize}
בנקודות $x=2\pm\sqrt2$ הפונקציה רציפה ו-$f''$ מחליפה סימן, ולכן אלו \textbf{נקודות פיתול}:
\[ \left(2-\sqrt2,\ (2-\sqrt2)^2e^{-(2-\sqrt2)}\right)=\left(2-\sqrt2,\ (6-4\sqrt2)e^{\sqrt2-2}\right), \]
\[ \left(2+\sqrt2,\ (6+4\sqrt2)e^{-2-\sqrt2}\right). \]
\end{solution}
